\clearpage
\section{Implementation and Results}
\label{sec:implementation}

The code is written in Python with NumPy only (no SciPy, no sparse library).
It is in our repository~\cite{repo}. The files are listed in Table~\ref{tab:files}.

\begin{table}[H]\centering
\begin{tabular}{ll}
\toprule
\textbf{File} & \textbf{Content} \\
\midrule
\texttt{toolkit.py} & the matrices, the sparse storage, the power method, the file reader \\
\texttt{webs.py} & every small web used in this report, written as a dictionary \\
\texttt{run\_paper.py} & tests against the numbers printed in the paper \\
\texttt{hollins\_main.py} & PageRank on the Hollins dataset \\
\texttt{ex01.py}, \dots, \texttt{ex17.py} & one script per computational exercise \\
\bottomrule
\end{tabular}
\caption{Files of the project}
\label{tab:files}
\end{table}

\subsection{Data structures}

\subsubsection*{A web is a dictionary}

A web is stored as a Python dictionary \verb|{page: [pages it links to]}| with pages numbered $1,\dots,n$ as in the paper.
For example Figure~2.1 is
\begin{verbatim}
FIG_2_1 = {1: [2, 3, 4], 2: [3, 4], 3: [1], 4: [1, 3]}
\end{verbatim}
A page with an empty list is a dangling node.

\subsubsection*{Dense matrix for small webs}

The function \texttt{build\_A} fills an $n \times n$ NumPy array with $A_{ij} = 1/n_j$.
We use it only for the small webs, to print $\A$ and $\M$ and to compute all eigenvalues with \texttt{numpy.linalg.eig} for cross-checking (for example to find $\dim\Vone{\A}$).
The dense matrix is \emph{not} used by the solver.

\subsubsection*{Sparse matrix (CSR) for the dataset}

For Hollins a dense $\A$ would need $6012^2 \cdot 8$ bytes $\approx 289$~MB, but only $23\,875$ entries are non-zero.
We store $\A$ in the \emph{compressed sparse row} format with three arrays, built by hand in \texttt{build\_csr}:
\begin{itemize}
  \item \texttt{AA}: the non-zero values of $\A$, read row by row;
  \item \texttt{JA}: the column index of each value in \texttt{AA};
  \item \texttt{IA}: for each row $i$, the position in \texttt{AA} where row $i$ starts (length $n+1$).
\end{itemize}
The non-zeros of row $i$ are \texttt{AA[IA[i]:IA[i+1]]}.
The three arrays for Hollins take $0.43$~MB, about 670 times less than the dense matrix.
The function \texttt{csr\_matvec} computes $\mathbf{y} = \A\vx$ row by row:
\[
  y_i = \sum_{k=\texttt{IA}[i]}^{\texttt{IA}[i+1]-1} \texttt{AA}[k]\; x_{\texttt{JA}[k]} .
\]
We tested that this product equals the dense product \texttt{A @ v} on a random vector (difference $0$).

\subsection{The solver}

The only solver of the project is the function \texttt{pagerank(links, n, m, tol)}.
It starts from $\vx_0 = (1/n,\dots,1/n)^T$ and repeats three steps until $\|\vx_{k+1}-\vx_k\|_1 < \texttt{tol}$:
\begin{enumerate}[label=\textbf{Step \Alph*.},leftmargin=*]
  \item $\mathbf{y} = \A\vx_k$ with the CSR product.
  \item $\mathbf{y} \leftarrow \mathbf{y} + \frac{1}{n}\sum_{j\ \text{dangling}} (\vx_k)_j$: the mass of the dangling pages is given back, equally to all pages.
  \item $\vx_{k+1} = (1-m)\mathbf{y} + m\vs$: this is formula (3.2).
\end{enumerate}
After Step~C the sum of $\vx_{k+1}$ is already $1$ in exact arithmetic; we divide by the sum once more only to remove rounding errors.
The function returns the vector $\vx$, the number of iterations $k$ and the last residual $\|\vx_{k+1}-\vx_k\|_1 = \|\M\vx_k - \vx_k\|_1$.
If the residual is not below \texttt{tol}, the method did not converge.

With $m = 0$ the same function iterates $\vx_{k+1} = \A\vx_k$, so we also use it to rank a web with formula (2.1) (Exercises~1 and~12).

\subsubsection*{Reading the dataset}

The function \texttt{load\_dat} reads \texttt{hollins.dat}: the first line gives $n$ and the number of links, then $n$ lines give the URL of each page, then one link per line as ``source target''.
Following the paper, a link from a page to itself would be ignored, and a repeated link would be counted once.
In the Hollins file there are no self-links and no repeated links.

\subsection{Tests on the two webs of the paper}

Before using the dataset we checked every number that the paper prints (script \texttt{run\_paper.py}).
Table~\ref{tab:paper-checks} shows that all of them agree.
For each value the allowed difference is one unit of the last printed decimal of the paper.

\begin{table}[H]\centering
\begin{tabular}{llcc}
\toprule
\textbf{Web} & \textbf{Quantity checked} & \textbf{Paper} & \textbf{Max.\ difference} \\
\midrule
Fig.~2.1 & eigenvector of $\A$ $= (12,4,9,6)/31$ & p.~3 & $1.7\cdot10^{-16}$ \\
Fig.~2.1 & $\dim\Vone{\A} = 1$ & p.~3 & exact \\
Fig.~2.1 & matrix $\M$ (four decimals) & Example 1 & $3.3\cdot10^{-6}$ \\
Fig.~2.1 & scores $0.368,\,0.142,\,0.288,\,0.202$ & Example 1 & $1.9\cdot10^{-4}$ \\
Fig.~2.2 & $\dim\Vone{\A} = 2$ & p.~4 & exact \\
Fig.~2.2 & matrix $\M$ of formula (3.3) & p.~6 & $5.6\cdot10^{-17}$ \\
Fig.~2.2 & scores $0.2,\,0.2,\,0.285,\,0.285,\,0.03$ & Example 2 & $5.6\cdot10^{-17}$ \\
\bottomrule
\end{tabular}
\caption{Checks of our code against the values printed in the paper}
\label{tab:paper-checks}
\end{table}


The ranking for Figure~2.1 with $\M$ is
\[
  \vx \approx [0.368151,\ 0.141809,\ 0.287962,\ 0.202078]
\]
\sumcheck{1.000}
and for Figure~2.2 it is exactly $[0.2,\ 0.2,\ 0.285,\ 0.285,\ 0.03]$, as in Example~2 of the paper.
The order of Figure~2.1 is $1 > 3 > 4 > 2$, the same order as with $\A$.

We also tested the dangling-node fix on a small web (Figure~2.1 without the link $3 \to 1$, so page~3 is dangling): the solver and a dense eigenvector computation on the corrected matrix give the same vector (difference $4\cdot10^{-14}$).

\subsection{Results on the Hollins dataset}
\label{sec:hollins}

Table~\ref{tab:hollins} summarises the run of \texttt{hollins\_main.py} with $m = 0.15$ and $\texttt{tol} = 10^{-12}$.

\begin{table}[H]\centering
\begin{tabular}{lr}
\toprule
\textbf{Quantity} & \textbf{Value} \\
\midrule
pages $n$ & 6012 \\
links (non-zeros of $\A$) & 23\,875 \\
fill of $\A$ & 0.066\% \\
dangling pages & 3189 (53\%) \\
memory, CSR arrays & 0.43 MB \\
memory, dense matrix & 289 MB \\
iterations ($m=0.15$, $\texttt{tol}=10^{-12}$) & 138 \\
final residual $\|\M\vx_k-\vx_k\|_1$ & $9.7\cdot10^{-13}$ \\
time (build CSR + iterations) & 1.0 s \\
\bottomrule
\end{tabular}
\caption{PageRank on the Hollins dataset}
\label{tab:hollins}
\end{table}


The most important page is the home page of the university, with a score about 120 times larger than the uniform score $1/n = 0.000166$.
Table~\ref{tab:hollins-top} lists the ten best pages.
The pages of the admissions area are high because many pages of the site link to them.

\begin{table}[H]\centering
\small
\begin{tabular}{crll}
\toprule
\textbf{Rank} & \textbf{Page} & \textbf{Score} & \textbf{URL} \\
\midrule
1 & 2 & 0.017669 & \texttt{hollins.edu/} \\
2 & 37 & 0.010238 & \texttt{hollins.edu/admissions/visit/visit.htm} \\
3 & 38 & 0.009457 & \texttt{hollins.edu/about/about\_tour.htm} \\
4 & 61 & 0.009069 & \texttt{hollins.edu/htdig/index.html} \\
5 & 52 & 0.008857 & \texttt{hollins.edu/admissions/info-request/info-request.cfm} \\
6 & 4023 & 0.008121 & \texttt{www1.hollins.edu/faculty/saloweyca/clas\%20395/Sculpture/sld001.htm} \\
7 & 43 & 0.007805 & \texttt{hollins.edu/admissions/apply/apply.htm} \\
8 & 27 & 0.006934 & \texttt{hollins.edu/admissions/admissions.htm} \\
9 & 3227 & 0.006924 & \texttt{www1.hollins.edu/faculty/saloweyca/clas\%20395/Sculpture/index.htm} \\
10 & 5254 & 0.006319 & \texttt{www1.hollins.edu/faculty/saloweyca/clas\%20395/Sculpture/tsld001.htm} \\
\bottomrule
\end{tabular}
\caption{The ten highest scores on Hollins ($m = 0.15$); the prefix \texttt{http://www.} is omitted}
\label{tab:hollins-top}
\end{table}




Every score is at least $m/n = 2.5\cdot10^{-5}$, as formula (3.2) predicts (the smallest score is $5.8\cdot10^{-5}$).
